\begin{abstract}
本文构造并审计一条闭合的形式推理链：以六维周期晶格的切投影（cut-and-project）方案为本体模型，以上位空间内部相位参数化“观察者”，以环面旋转的一维因子作为有限资源下的可实现采样机制，以二分割读出诱导的机械词（Sturmian）刻画最小复杂度二值信息流，并以 Zeckendorf 规范形定义稳定化核，将原始读出压缩至黄金均值合法语言。由此得到：时间可被形式化为稳定化核的预像纤维在观测历史下的逐步细化；其信息率由度量熵控制，并在遍历条件下呈现与观察者相位无关的统计律。本文进一步给出节点—纤维入射图的闭包动力学与可复现的数值协议；并在附录中提出“可见性投影”与“有效维数”框架，用以讨论非整数维的资源维度与黄金比标度，以及符号动力系统 $\zeta$ 生成函数与算术结构之间的可检验接口。
\end{abstract}

\begin{sloppypar}
\textbf{关键词：} 切投影；六维超空间；观察者参数；环面旋转；Sturmian 机械词；Zeckendorf 表示；稳定化核；预像纤维；度量熵；遍历性；统计全息
\end{sloppypar}

\tableofcontents
\newpage

